\documentclass[12pt]{article}
\usepackage[utf8]{inputenc}
\usepackage[T1]{fontenc}
\usepackage{hyperref}
\title{Santo Domingo Albarradas Zapotec Language Resources for K-12 Curriculum Development}
\author{Compiled by Claude Sonnet 5.5 (claude-sonnet-5-5) for Andr\'as Kornai}
\date{}
\begin{document}
\maketitle

\section{Introduction}
Santo Domingo Albarradas Zapotec (ISO 639-3 zas; Glottocode sant1450), also called Albarradas Zapotec, is a Central Zapotec language of Oaxaca, Mexico, spoken in Santa Mar\'ia, Santo Domingo and San Miguel Albarradas. Wikipedia cites about 5,500 speakers from the 1980 census, so the figure is old. The neighbouring varieties of Santa Catarina and San Antonio Albarradas are reported not to be mutually intelligible. Glottolog describes the language as shifting, with use growing in Santo Domingo and declining in Santa Mar\'ia. Wikitongues gives the endonym Dihidx Bily\'ahab.

\section{Primary Educational Resources}
No primers, textbooks or school programmes were found in the search results.

\section{Linguistic and Reference Materials}
\begin{itemize}
\item Wikipedia, Albarradas Zapotec (overview and classification). \url{https://en.wikipedia.org/wiki/Albarradas_Zapotec}
\item Glottolog entry for Santo Domingo Albarradas Zapotec. \url{https://glottolog.org/resource/languoid/id/sant1450}
\item Endangered Languages Project entry, with a linked bibliography. \url{https://endangeredlanguages.com/lang/8792}
\item Verbix endangered-languages page. \url{https://docs.verbix.com/EndangeredLanguages/AlbarradasZapotec}
\end{itemize}
No dictionary or grammar of this variety was confirmed.

\section{Religious and Community Resources}
Global Recordings Network reports Scripture and evangelistic audio, including Words of Life Bible stories and Faith Comes By Hearing New Testament recordings: \url{https://globalrecordings.net/en/language/zas}

\section{Digital and Audio Resources}
\begin{itemize}
\item Global Recordings Network page for zas (audio). \url{https://globalrecordings.net/en/language/zas}
\item Wikitongues page for the language. \url{https://wikitongues.org/languages/zas/}
\end{itemize}

\section{Curriculum Development Context}
Only reference-level and religious audio materials were found. A curriculum would need new materials developed with the communities, and the Albarradas varieties should be kept distinct from each other.

\end{document}